\chapter{Supervision Interface}

\section{Design Principles}

The web interface serves two distinct user profiles with different requirements:

\paragraph{Panel operator.}
An operator standing at the process, monitoring temperatures and flow rates, occasionally
adjusting setpoints. The SIMATIC Unified 7\textquotedblright{} touchscreen imposes a
resolution of 800~$\times$~480 pixels and a typical viewing distance of 0.5--1.5\,m.
Key requirements: at-a-glance process state, large readable values, touch-friendly controls
(minimum 44~$\times$~44\,px tap targets), no keyboard, no scroll.

\paragraph{Process engineer.}
An engineer commissioning the cascade controller or analysing step test results on a PC.
Key requirements: detailed process data, parameter forms, live trend charts, CSV export.
No screen size constraint.

These two profiles led to designing separate primary interfaces:
the \textbf{Synoptique} (\code{/schema}) for the panel and
the \textbf{Cascade} (\code{/cascade}) for the engineer.
Both pages connect to the same Socket.IO backend and update every 3\,s.

Common visual design decisions across all pages:
\begin{itemize}[itemsep=3pt]
  \item Dark header bar with product name and data source badge
    (\textcolor{cmtgreen}{\texttt{OPENPIPE}} / \textcolor{orange!80!black}{\texttt{SIMULATION}})
  \item Monospaced font for numerical values (prevents layout shifts when digit count changes)
  \item Colour-coded badges for discrete states (valve position, controller mode)
  \item Bootstrap~5 for responsive layout; Chart.js~4 for live charts
\end{itemize}

\section{Synoptique (\texttt{/schema}) --- Panel Display}

The Synoptique is a full-screen SVG-based P\&ID diagram of the installation. The diagram
shows the furnace on the left, HE-001 in the centre, the hot circuit flowing
left-to-right at the top, and the cold circuit flowing right-to-left at the bottom.

\textbf{Animated flow lines.} CSS-animated dashed strokes create a visual flow
indication: orange for the hot circuit (F1), blue for the cold circuit (F2).
The animation runs continuously (visual only; no speed modulation with flow rate).

\textbf{Live value overlays.} Ten \code{<foreignObject>} HTML islands are positioned over
the corresponding points of the P\&ID diagram. Each card shows the tag name, current value,
and unit. All ten values update on each \code{process\_tags} Socket.IO event.

\textbf{Furnace power indicator.} A coloured dot (green~=~ON, red~=~OFF) driven by the
\code{power} Boolean tag.

\textbf{Setpoint controls.} Two pairs of $+$/$-$ buttons adjust $F_{1,SP}$ and
$F_{2,SP}$ in \SI{0.1}{\metre\cubed\per\hour} steps. Each press sends a REST call to
\code{POST /api/process/write} with the tag name and new value.

\section{Cascade Page (\texttt{/cascade}) --- Engineering Tool}

\subsection{Identification Tab}

The identification tab implements the full workflow from step test configuration to IMC
gain application.

\textbf{Step test form.}
Three inputs: Base~$F_{1,SP}$ (auto-filled from the current PLC value when the tab loads),
Amplitude (\si{\metre\cubed\per\hour}), and Duration (\si{\second}). The \textit{Start Test}
button calls \code{POST /api/identification/start} and is disabled when mode $\neq$
MANUAL.

\textbf{Progress display.}
A two-phase progress bar tracks the current state: PRE-STEP (grey, 15\,s baseline) and
STEP (blue, configurable duration).

\textbf{Live chart.}
A Chart.js dual-axis line chart accumulates data from the start of the test.
Left~Y-axis: $T_{in1,HE1}$ (red) and $T_{out2,HE1}$ (orange dashed).
Right~Y-axis: $F_{1,SP}$ step (blue stepped line). New data points are appended on each
\code{ident\_update} Socket.IO event using \code{chart.update('none')} for no-animation
streaming.

\textbf{Results display.}
When the test completes: three KPI cards for $K$, $\tau$, $\theta$; an IMC tuning table
with three rows (Aggressive / Normal / Conservative); and an \textit{Apply~\&~Switch} button
per row that posts the gains to \code{POST /api/cascade/params} and activates the Cascade
Control tab.

\subsection{Cascade Control Tab}

\textbf{Block diagram.}
A CSS-styled visual representation of the cascade structure:
$T_{out2,SP} \rightarrow [P] \rightarrow T_{in1,SP} \rightarrow [\text{PI+P}]
\rightarrow F_{1,SP}$.
Each block displays its current gain values.

\textbf{Parameter forms.}
Inner loop ($K_p$, $T_i$) and outer loop ($K_{p,ext}$, $T_{in1,SP}^{\text{base}}$)
inputs. Changes apply via \code{POST /api/cascade/params} and take effect at the next
poll cycle.

\textbf{Mode buttons.}
Three large buttons for MANUAL, AUTO-INNER, and AUTO-FULL. The active mode is highlighted.
A confirmation dialog is shown before switching from any AUTO mode back to MANUAL.

\textbf{Status grid.}
Live display of inner loop state ($T_{in1,SP}$, $T_{in1}$, error) and outer loop state
($T_{out2,SP}$, $T_{out2}$, error), plus current $F_{1,SP}$ output.
All values update on each \code{cascade\_data} Socket.IO event.

\section{Real-Time Performance}

\textbf{Socket.IO latency.}
On the laboratory LAN (192.168.1.x), end-to-end latency from PLC poll to browser display
is dominated by the 3-second poll period. The WebSocket push itself adds $<$\,10\,ms.
This is adequate for a thermal process with time constants of tens of seconds.

\textbf{Chart.js rendering.}
The identification chart accumulates up to 200 data points (10\,minutes at 3\,s/sample)
without noticeable lag, using \code{animation: false} and the
\code{chart.update('none')} mode to append new points without re-rendering the full chart.

\textbf{Panel browser compatibility.}
The application has been tested in the Chromium-based browser embedded in the SIMATIC
Unified 7\textquotedblright{} panel runtime. Bootstrap~5 and Chart.js~4 are both
compatible. Socket.IO WebSocket connections work over the panel's LAN interface.

\section{Deployment Evolution}
\label{sec:deploy_evolution}

The supervision application went through three successive deployment environments before
reaching its target configuration. Each stage revealed new constraints that shaped the
final architecture.

\subsection{Stage 1 --- Raspberry Pi 4 (Standalone Validation)}

The first deployment targeted a Raspberry Pi~4 (4\,GB RAM, \code{192.168.1.38}) connected
to the laboratory switch alongside the S7-1500 and the SIMATIC Unified panel.
The objective was to validate the entire supervision stack in a \emph{standalone}
configuration, with no dependency on Siemens-proprietary infrastructure
(no Industrial Edge licence, no WinCC Unified runtime involvement).

On this platform, the Flask application was containerised with Docker and communicated
with the PLC exclusively via Snap7 over TCP/IP. This stage confirmed:
\begin{itemize}[leftmargin=*, itemsep=3pt]
  \item \textbf{Snap7 connectivity:} DB1 reads and writes functional over the laboratory
    LAN at \code{192.168.1.100}.
  \item \textbf{Real-time pipeline:} Socket.IO push events reaching connected browsers
    every 3\,s with no observed latency on the LAN.
  \item \textbf{Docker on ARM:} the multi-platform image built on an x86 PC
    ran correctly on the Pi's ARM Cortex-A72 processor without modification.
  \item \textbf{Browser accessibility:} the panel's Chromium-based browser successfully
    rendered the Synoptique at \code{http://192.168.1.38:5000}.
\end{itemize}

The Pi deployment also revealed the boundary of the standalone approach:
the Open Pipe Unix socket (\code{/tmp/siemens/automation}) is created exclusively by
the WinCC Unified runtime on the panel hardware; it is not reachable from an
external device. On the Pi, all Open Pipe calls fell back to the simulation stub,
meaning WinCC tag write operations (setpoints, controller mode) could not reach the PLC
through this path.

\textit{Outcome:} Standalone Snap7 supervision fully validated on ARM hardware;
Open Pipe integration confirmed to require local execution on the panel itself,
motivating the Industrial Edge deployment target.

\subsection{Stage 2 --- Development PC with Docker (Staging Environment)}

The application was containerised and deployed on the development PC
(\code{192.168.1.149}), which runs on the same LAN as the PLC and panel.
Docker provided three key advantages over the Pi prototype:
\begin{itemize}[leftmargin=*, itemsep=3pt]
  \item \textbf{Reproducibility:} the \code{docker-compose.yml} file captures the entire
    runtime environment; the image built on the PC is binary-identical to the one that
    will run on the panel.
  \item \textbf{Hot-reload development:} the \code{./app} directory is bind-mounted into
    the container; Python file changes take effect without rebuilding the image.
  \item \textbf{Panel accessibility:} a Windows Firewall inbound rule was added to allow
    TCP connections on port~5000 from the local subnet, making the interface immediately
    accessible from the SIMATIC Unified panel browser at \code{http://192.168.1.149:5000}.
\end{itemize}

At this stage, Open Pipe write access was simulated via the Python stub in
\code{openpipe\_client.py}: the socket path \code{/tmp/siemens/automation} does not
exist on the PC, so all Open Pipe calls fall back gracefully without blocking the
application. Snap7 DB1 reads and writes operate normally over TCP.

\textit{Outcome:} Full application functionality validated in a production-equivalent
Docker environment; identified the \code{/tmp/siemens/automation} socket as the
final integration dependency.

\subsection{Stage 3 --- Industrial Edge on SIMATIC Unified Panel (Target)}

The final deployment target places the Docker container natively on the
SIMATIC Unified 7\textquotedblright{} panel via Industrial Edge. In this configuration:
\begin{itemize}[leftmargin=*, itemsep=3pt]
  \item The \code{/tmp/siemens/automation} socket is accessible inside the container
    (volume mount in \code{docker-compose.edge.yml}), enabling full Open Pipe read/write.
  \item The panel browser accesses the application at \code{http://localhost:30500}
    (same-device loopback), eliminating any dependency on a separate PC.
  \item The installation becomes fully self-contained: PLC, panel, and supervision
    application form a single unit with no external computer required.
\end{itemize}

\textbf{Current status:} The \texttt{.app} package has been exported from App Publisher
and validated. Deployment is pending Industrial Edge licence acquisition for the panel at
\code{192.168.1.200}.

\section{Deployment Summary}

\begin{table}[H]
  \centering
  \caption{Web application deployment environments.}
  \label{tab:deploy}
  \begin{tabular}{@{}L{4.5cm}L{5.5cm}L{3.5cm}@{}}
    \toprule
    \textbf{Environment} & \textbf{URL} & \textbf{Status} \\
    \midrule
    Development PC (localhost)
      & \code{http://localhost:5000}
      & Active \\
    PC $\rightarrow$ Panel (LAN)
      & \code{http://192.168.1.149:5000}
      & Active (firewall rule added) \\
    Panel Edge (target)
      & \code{http://192.168.1.200:30500}
      & Pending (licence required) \\
    \bottomrule
  \end{tabular}
\end{table}
